\documentclass[11pt,a4paper,fontset=windows]{ctexart}
\usepackage[margin=2.2cm]{geometry}
\usepackage{booktabs,array,longtable,amsmath}
\usepackage[hidelinks]{hyperref}
\setlength{\parskip}{4pt}
\title{BLIS 如何表示异构集群：现状、机制与做异构研究需要改的东西}
\author{ours 分支（上游 2ebef6a），2026-09-26}
\date{}
\begin{document}
\maketitle

\section{结论摘要}

\begin{enumerate}
\item \textbf{我们目前的仿真是同构集群。} 每台实例（instance，一个运行 vLLM 引擎的推理副本）都是 Qwen3-14B、H100、张量并行度 1（TP，一个副本占一块 GPU）、同样的 KV 块数、同样的步时长模型。系统规模 $n$ 就是 $n$ 台一模一样的单 GPU 服务器。
\item \textbf{BLIS 有"节点池"（node pools）机制来声明多种 GPU 各多少台，但它只改变三样东西：每台实例的 KV 容量、最大上下文长度、每小时成本。} 步时长（一步计算要多久）不会随 GPU 类型改变：代码里预留了"按 GPU 覆盖延迟系数"的字段，但命令行和策略文件都没有地方填它（上游 issue \#893 承认此事）。所以今天用节点池混合 H100 与 L40S，L40S 只是"内存小的 H100"。
\item \textbf{路由看不见 GPU 类型。} 路由快照里有 GPU 类型和成本字段，但所有评分器只读队列长度和 KV 利用率（比例）。一台 4 千块的 L40S 和一台 1.6 万块的 H100 在利用率 50\% 时被视为相同。
\item \textbf{PD 分离（prefill 与 decode 分池）是 BLIS 里唯一"按角色"的异构}，可按角色设 TP、KV 容量、上下文上限、延迟后端，但同样不能按角色设 GPU 速度。
\item \textbf{我们的共享队列模式在异构下能跑但有隐含偏向}：唤醒最低编号实例、fcfs-pool 不分实例、Mooncake 用 0 号实例的步时长模型、后加入的实例拿不到共享队列。
\item 做异构研究的最小改动是一处：在 \texttt{cmd/root.go} 里为每个节点池加载对应 GPU 的硬件校准并填入 \texttt{HWConfigByGPU}。此外要放弃 \texttt{--total-kv-blocks}（否则所有实例同容量），并决定唤醒规则与集合策略是否应感知实例速度。
\end{enumerate}

\section{我们现在跑的是什么}

\texttt{ours/sweep.py} 的基本参数：\texttt{--model qwen/qwen3-14b --hardware H100 --tp 1 --num-instances n --total-kv-blocks B}，加共享队列（默认）或 push 路由（B1 系列），不传 \texttt{--policy-config}。因此：没有节点池，没有放置，没有成本，没有拓扑；所有实例用同一套 H100 校准（trained-physics 后端，用 H100 的峰值算力 989.5 TFLOPS 和显存带宽 3.35 TB/s 做基函数，系数在 \texttt{defaults.yaml} 里只有一套，是在 H100 轨迹上拟合的）；每台实例 $B$ 块（15,909 或 6,500），由命令行显式给定，而不是按显存自动算。

\section{BLIS 里的异构机制}

\subsection{节点池：如何声明"每种 GPU 各几台"}

在 \texttt{--policy-config} 的 YAML 里写 \texttt{node\_pools}（类型 \texttt{NodePoolConfig}，\texttt{sim/cluster/infra\_config.go}）。字段：

\begin{center}\begin{tabular}{ll}\toprule
字段 & 含义 \\ \midrule
\texttt{name} & 池名，必须唯一 \\
\texttt{gpu\_type} & GPU 类型字符串，必须跨池唯一（成本与容量按它查） \\
\texttt{gpus\_per\_node} & 每节点 GPU 数（$\ge 1$） \\
\texttt{gpu\_memory\_gib} & 每块 GPU 显存，用于自动算 KV 容量 \\
\texttt{initial\_nodes} & 起始节点数（就是"服务器数"） \\
\texttt{min\_nodes}, \texttt{max\_nodes} & 上下界；\texttt{max\_nodes} 必须 $\ge$ \texttt{initial\_nodes}，所以必须写 \\
\texttt{provisioning\_delay} & 新节点上线延迟（在 \texttt{blis run} 里没有代码会新建节点，等于无用） \\
\texttt{cost\_per\_hour} & 每节点每小时成本 \\ \bottomrule
\end{tabular}\end{center}

一个最小例子（4 台 H100、8 台 L40S，每台一块 GPU，TP = 1）：
\begin{verbatim}
node_pools:
  - {name: h100, gpu_type: H100, gpus_per_node: 1, gpu_memory_gib: 80,
     initial_nodes: 4, max_nodes: 4, cost_per_hour: 4.0}
  - {name: l40s, gpu_type: L40S, gpus_per_node: 1, gpu_memory_gib: 48,
     initial_nodes: 8, max_nodes: 8, cost_per_hour: 1.2}
instance_lifecycle:
  warm_start_initial_instances: true
\end{verbatim}
运行时仍要 \texttt{--num-instances 12}：实例数来自命令行，不是池的总和。\texttt{gpu\_memory\_gib} 与 \texttt{hardware\_config.json} 里的 \texttt{MemoryGiB} 是两处独立的数，要手工保持一致。

\subsection{放置：哪台实例落到哪个池}

\texttt{PlaceInstance}（\texttt{sim/cluster/infra\_placement.go}）按实例编号 0, 1, 2, \ldots 依次放置，第一遍在"单节点首次适配"下按池的\textbf{声明顺序}、节点编号、空闲 GPU 顺序找位置；第二遍只在 TP 大于每节点 GPU 数时跨整节点放置。含义：
\begin{itemize}
\item 先声明的池先被填满。上例里实例 0 到 3 落在 H100，4 到 11 落在 L40S。
\item 没有"把实例 $i$ 钉在池 $j$"的办法，唯一的控制是池的顺序和数量。
\item 实例数超过总容量时，多出的实例进入 pending，只打一条警告，此后永远不会被放置（\texttt{NodeReadyEvent} 在 \texttt{blis run} 里没有生产代码触发），相当于被静默丢弃。
\end{itemize}

\subsection{放置之后每台实例得到什么}

三处放置代码（启动、延迟放置、自动扩容）做同样的几步：
\begin{enumerate}
\item \textbf{GPU 标签}：\texttt{simCfg.GPU} 设为池的 \texttt{gpu\_type}。只是标签。
\item \textbf{延迟系数（不生效）}：代码为 \texttt{if hc, ok := config.HWConfigByGPU[type]; ok \{ simCfg.HWConfig = hc \}}。\texttt{HWConfigByGPU} 在 \texttt{DeploymentConfig} 里有定义，但 \texttt{cmd/} 从未赋值，策略文件也没有对应的键（\texttt{sim/cluster/network\_topology.go} 的注释明说：mixed-gpu\_type fleet currently shares one calibration）。因此每台实例都用 \texttt{--hardware} 那一种 GPU 的校准算步时长。
\item \textbf{KV 容量}（\texttt{applyPerInstanceKVCapacity}，\texttt{sim/cluster/kv\_autocalc.go}）：仅当没有显式传 \texttt{--total-kv-blocks} 时，用池的 \texttt{gpu\_memory\_gib} 重算：
\[\text{blocks} = \frac{\text{显存}\times\text{利用率}\times TP - (\text{权重} + \text{激活 5.5 GiB} + \text{非 torch 开销})}{16\ \text{tokens}\times\text{每 token KV 字节}}\]
用同一 fork 二进制实测 Qwen3-14B、TP 1、利用率 0.9：H100 与 A100-80 都是 15,909 块，L40S 4,113 块。显式给了 \texttt{--total-kv-blocks} 则所有实例同容量。
\item \textbf{上下文上限}：若 \texttt{max-model-len} 超过该实例的 KV 容量，就按容量封顶（小 GPU 的实例会拒绝更长的请求）。
\item \textbf{成本}：\texttt{InstanceCostPerHour} 等于池的每小时成本乘以跨越的节点数。它不进任何输出指标，只供自动扩缩容和路由快照使用。
\end{enumerate}

\subsection{步时长如何依赖 GPU（trained-physics 后端）}

步时长 = 一组 roofline 基函数（按 token 数、上下文长度、层数算出的计算量和访存量，分别除以 GPU 的峰值算力 \texttt{TFlopsPeak} 和显存带宽 \texttt{BwPeakTBs}）乘以一套拟合系数 $\alpha,\beta$。GPU 类型只通过这两个峰值（以及权重为 1 字节时的 \texttt{TFlopsFP8}，多节点时的互连带宽比）进入。随附的 \texttt{hardware\_config.json} 只有四个条目：

\begin{center}\begin{tabular}{lrrr}\toprule
GPU & 峰值 TFLOPS (bf16) & 显存带宽 TB/s & 显存 GiB \\ \midrule
H100 & 989.5 & 3.35 & 80 \\
A100-SXM / A100-80 & 312 & 2.04 & 80 \\
L40S & 362 & 0.86 & 48 \\ \bottomrule
\end{tabular}\end{center}

系数 $\alpha,\beta$ 全局只有一套，在 H100 上拟合；\texttt{defaults.yaml} 里 30 个模型条目全部标 H100。所以 A100 或 L40S 的绝对延迟是"用 H100 系数按峰值等比缩放"的结果，没有校准过。作为研究里的"快机器与慢机器"是可用的，但不要把它们的绝对数字当成 A100 或 L40S 的实测。

\subsection{路由能否区分快慢实例}

路由快照 \texttt{RoutingSnapshot} 携带 \texttt{GPUType}、\texttt{TPDegree}、\texttt{CostPerHour}、\texttt{TotalKvCapacityTokens}、\texttt{KvTokensInUse}、\texttt{MaxBatchSize}（在 \texttt{cluster\_event.go} 填入）。但没有任何评分器读它们：

\begin{center}\begin{tabular}{ll}\toprule
评分器 & 读取的量 \\ \midrule
queue-depth, load-aware & 队列长度 \\
kv-utilization & $1 - $ 已用块 / 总块（比例） \\
active-requests, running-requests & 在途数、批大小 \\
vllm-dp & $4\times$队列长度 $+$ 批大小 \\ \bottomrule
\end{tabular}\end{center}

快照里的 \texttt{TTFT}、\texttt{ITL}、\texttt{DispatchRate} 字段从未被生产代码赋值，恒为 0。因此今天的路由策略只能通过积压间接感知快慢；要做异构路由需要新的评分器（读 GPU 类型、绝对容量或步时长）。

\subsection{PD 分离：按角色的异构}

\texttt{--prefill-instances p --decode-instances d}（另有共享角色）按编号分配角色：先 prefill、再 decode、再共享。每角色可覆盖 \texttt{tp}、\texttt{hardware}、\texttt{latency-model}、\texttt{max-model-len}、MoE 通信后端（\texttt{ResolvePoolConfig}，\texttt{sim/cluster/resolve.go}）。其中 \texttt{--prefill-hardware L40S} 同样只改标签和该角色的 KV 自动计算，不改延迟校准；能真正进入步时长的角色差异只有 TP 和延迟后端。KV 传输成本 = 基础延迟 0.05 ms + 传输字节 / 25 GB/s（可开争用：带宽按并发传输数均分），与 GPU 类型和放置无关。注意 PD 与节点池同时用时放置不分角色（角色按编号，池按顺序）。

\subsection{自动扩缩容与成本}

生产接线是 \texttt{DefaultCollector} $\to$ \texttt{V2SaturationAnalyzer} $\to$ \texttt{UnlimitedEngine} $\to$ \texttt{DirectActuator}。分析器按 (GPU 类型, TP) 分组比较"KV 供给 = 总 KV token $\times$ 阈值"与"KV 需求 = 在用 token + 队列长度 $\times$ 平均输入"；引擎扩最便宜的变体、缩最贵的。它只在现有节点的空闲 GPU 上放实例，从不新建节点；没有延迟或速度模型。成本没有任何输出指标，要自己按每小时成本之和乘以时长来算。

\subsection{限制}

\begin{center}\begin{tabular}{lll}\toprule
组合 & 结果 & 位置 \\ \midrule
\texttt{blis replay} + 节点池 & 致命错误 & \texttt{cmd/replay.go} \\
\texttt{blis observe} & 无 \texttt{--policy-config} 参数 & \\
MoE \texttt{--dp > 1} + 节点池 & 致命错误 & \texttt{cmd/root.go} \\
角色 TP $\ne$ 全局 TP + 节点池 & panic & \texttt{sim/cluster/cluster.go} \\
重复 \texttt{gpu\_type} 或池名 & panic & \texttt{infra\_placement.go} \\ \bottomrule
\end{tabular}\end{center}

\section{我们的共享队列模式在异构下的问题}

\begin{enumerate}
\item \textbf{命令行检查}只拒绝 \texttt{--flow-control} 和 prefill/decode 分池，不检查 \texttt{--prefill-decode-instances}；节点池可以直接与共享队列同用。
\item \textbf{唤醒规则}：到达和完成都唤醒"最低编号的空闲实例"，而放置先填先声明的池，所以 YAML 里池的顺序变成了 GPU 类型偏向：轻载时工作总是先给第一个池。
\item \textbf{fcfs-pool 不知道是谁在取}：\texttt{BatchContext} 没有实例身份或速度字段。快实例更频繁到达步边界，隐含地多拉工作，这是"自然的"，但任何想按速度分配的集合策略都需要新增字段。
\item \textbf{按实例的上下文上限丢弃}：\texttt{AdoptSharedRequest} 若请求超过本机的 \texttt{MaxModelLen} 或 KV 容量就丢弃。小 GPU 实例会丢掉大 GPU 本可服务的请求，且与自检报告 A2 的"先分配再取件"问题叠加。
\item \textbf{Mooncake 用 0 号实例的步时长模型}（\texttt{cluster.go}，注释写着 all instances share one config），异构下预测是错的。alpha 延迟取 0 号实例无妨，因为该延迟是全局常数。
\item \textbf{后加入的实例拿不到共享队列}：\texttt{SetSharedQueue} 只在启动时调用，自动扩容或延迟上线的实例永远不会从共享队列取件。异构研究若只用启动时的固定实例集则无影响。
\item 实例加载延迟非零时，\texttt{InstanceLoadedEvent} 不唤醒共享队列；用 \texttt{warm\_start\_initial\_instances: true} 或零延迟。
\end{enumerate}

\section{做异构研究需要改什么}

\begin{enumerate}
\item \textbf{按 GPU 的步时长（唯一的硬阻塞）}：在 \texttt{cmd/root.go} 组装 \texttt{DeploymentConfig} 处，对每个节点池调用 \texttt{latency.GetHWConfig(hwConfigPath, pool.gpu\_type)} 填入 \texttt{HWConfigByGPU}。放置代码三处已经会读它。约 20 行加一个测试。这一步之后，L40S 实例才会真的慢（带宽 0.86 对 3.35 TB/s，解码步长约 3 到 4 倍）。
\item \textbf{KV 容量}：去掉 \texttt{--total-kv-blocks} 让自动计算生效（并按池设 \texttt{gpu\_memory\_gib}）；或者新增"按池显式块数"（目前没有）。注意我们现在的 6,500 块远低于 H100 的物理容量 15,909，若要保持"内存受限"的研究设定，需要按池给显式块数，这一项也要加。
\item \textbf{放置控制}：按想要的顺序声明池，\texttt{--num-instances} 恰好等于总容量；多余实例被静默丢弃。
\item \textbf{共享队列}：唤醒规则改为轮转或按速度；给 \texttt{SetPolicy} 的上下文加实例身份、容量、步时长；\texttt{AdoptSharedRequest} 的按机丢弃要改成"放回队列给别的实例"或在入队时按最大实例检查；Mooncake 改为每实例的步时长函数；\texttt{addLiveInstance} 补 \texttt{SetSharedQueue}。
\item \textbf{路由基线}：现有评分器都是比例或计数，需要一个感知容量或速度的评分器作为 push 模式的异构基线。
\item \textbf{校准诚实}：非 H100 的绝对延迟未校准，报告里把它们写成"慢 GPU 变体（峰值等比缩放）"。
\item \textbf{成本}：没有成本输出，要在 \texttt{analyze.py} 里按池成本乘以窗口长度补一列，目标函数里才有"每小时成本"项。
\item \textbf{回放不可用}：节点池只在 \texttt{blis run} 里有效。
\end{enumerate}

\section{建议（供你决定，不会自行执行）}

若下一步研究"快慢混合的 GPU 池里如何分配请求"，最小可行形态是：两个池（H100 与 L40S，或 H100 与 A100-80 作为"同显存但慢 3 倍"的干净对照），启动时固定实例数，先做第 1 项改动并用单实例验证 L40S 的步时长确实变慢，再决定共享队列的唤醒规则和集合策略如何感知实例速度。这里的每一项选择（池的比例、显式块数、唤醒规则、集合策略）都是研究判断，需要你来定。

\end{document}
